\documentclass[10pt]{article}
\usepackage[margin=2.2cm]{geometry}
\usepackage{longtable}
\usepackage{booktabs}
\usepackage{amsmath}
\title{Supplementary Material: Prevalence Survey Coding Sheet for RUL Evaluation Protocols}
\author{Zhongkuan Ma}
\date{Draft for author verification}

\begin{document}
\maketitle

\section*{Coding Rules}
The survey covers 20 bearing RUL prediction papers published between 2019 and 2024. Each paper is coded by reading the methods and experiments sections.

\begin{enumerate}
\item \textbf{Piecewise 80\%}: the paper uses a piecewise linear RUL label with an 80\% plateau (or reports it as the standard piecewise RUL convention).
\item \textbf{Random split}: the paper randomly assigns sliding-window samples to train/validation/test without per-bearing holdout.
\item \textbf{Combined}: both piecewise labels and random time-step splitting are used in the reported main evaluation.
\item \textbf{Constant baseline}: the paper reports a constant predictor or another trivial baseline for calibration.
\end{enumerate}

\section*{Aggregate Counts}
\begin{center}
\begin{tabular}{lc}
\toprule
Criterion & Count \\
\midrule
Piecewise 80\% label & 18/20 (90\%) \\
Random time-step split & 16/20 (80\%) \\
Combined piecewise + random split & 14/20 (70\%) \\
Constant predictor baseline & 0/20 (0\%) \\
\bottomrule
\end{tabular}
\end{center}

\section*{Coding Sheet}
This draft lists the papers identified from the local reference set. Rows P01--P09 have paper metadata; rows P10--P20 require the author to insert the paper identifiers from the original 20-paper survey. All coding decisions require full-text verification before journal submission.

\small\setlength{\tabcolsep}{3pt}\begin{longtable}{p{0.8cm} p{2.8cm} p{1.2cm} p{3.0cm} p{1.2cm} p{1.2cm} p{1.2cm} p{2.2cm}}
\toprule
ID & First Author & Year & Journal & Piecewise & Random & Baseline & Evidence \\
\midrule
\endfirsthead
\toprule
ID & First Author & Year & Journal & Piecewise & Random & Baseline & Evidence \\
\midrule
\endhead
P01 & Wang & 2019 & IEEE Access & PENDING & PENDING & PENDING & Full text required \\
P02 & Li & 2020 & IEEE TII & PENDING & PENDING & PENDING & Full text required \\
P03 & Zhang & 2020 & MSSP & PENDING & PENDING & PENDING & Full text required \\
P04 & Chen & 2021 & Neurocomputing & PENDING & PENDING & PENDING & Full text required \\
P05 & Wang & 2021 & Measurement & PENDING & PENDING & PENDING & Full text required \\
P06 & Li & 2021 & IEEE Access & PENDING & PENDING & PENDING & Full text required \\
P07 & Mo & 2022 & IEEE TIM & PENDING & PENDING & PENDING & Full text required \\
P08 & Costa & 2022 & RESS & PENDING & PENDING & PENDING & Full text required \\
P09 & Li & 2022 & RESS & PENDING & PENDING & PENDING & Full text required \\
P10 & To be completed & & & PENDING & PENDING & PENDING & Full text required \\
P11 & To be completed & & & PENDING & PENDING & PENDING & Full text required \\
P12 & To be completed & & & PENDING & PENDING & PENDING & Full text required \\
P13 & To be completed & & & PENDING & PENDING & PENDING & Full text required \\
P14 & To be completed & & & PENDING & PENDING & PENDING & Full text required \\
P15 & To be completed & & & PENDING & PENDING & PENDING & Full text required \\
P16 & To be completed & & & PENDING & PENDING & PENDING & Full text required \\
P17 & To be completed & & & PENDING & PENDING & PENDING & Full text required \\
P18 & To be completed & & & PENDING & PENDING & PENDING & Full text required \\
P19 & To be completed & & & PENDING & PENDING & PENDING & Full text required \\
P20 & To be completed & & & PENDING & PENDING & PENDING & Full text required \\
\bottomrule
\end{longtable}

\section*{Status}
This is a draft for author verification. The aggregate counts used in the manuscript were produced in the earlier survey; the paper-level coding sheet must be completed and validated against full texts before submission. The manuscript should not claim that the coding sheet is final until all 20 rows are coded.

\end{document}
